% diversity_reduction_theory.tex
% Intended to be \input{} or \include{} into a main paper file.
% Assumes amsmath, amssymb, amsthm are loaded.
% Assumes theorem environments (theorem, proposition, definition, remark, proof) are defined.
%
% === INTEGRATION NOTES ===
%
% Citation keys used (match to your .bib entries):
%   \citet{brody2023expressivity}  -- Brody et al., "On the Expressivity Role of LayerNorm
%                                     in Transformers' Attention", ACL 2023
%   \citet{turntrout2023steering}  -- Turner et al., "Steering GPT-2-XL by Adding an
%                                     Activation Vector", LessWrong/Alignment Forum, May 2023
%   \citet{turntrout2023norms}     -- Turner & Ray, "Residual Stream Norms Grow Exponentially
%                                     over the Forward Pass", LessWrong/Alignment Forum, May 2023
%   \citet{dong2021attention}      -- Dong et al., "Attention is Not All You Need: Pure
%                                     Attention Loses Rank Doubly Exponentially with Depth",
%                                     ICML 2021
%   \citet{wu2024role}             -- Wu et al., "On the Role of Attention Masks and LayerNorm
%                                     in Transformers", NeurIPS 2024
%   \citet{sinii2025small}         -- Sinii et al., "Small Vectors, Big Effects: A Mechanistic
%                                     Study of RL-Induced Reasoning via Steering Vectors",
%                                     arXiv:2509.06608, 2025
%   \citet{zhang2025verbalized}    -- Zhang et al., "Verbalized Sampling: How to Mitigate Mode
%                                     Collapse and Unlock LLM Diversity", arXiv:2510.01171, 2025
%   \citet{omahony2024attributing} -- O'Mahony et al., "Attributing Mode Collapse in the
%                                     Fine-Tuning of LLMs", ICLR ME-FoMo Workshop 2024
%
% Required packages: amsmath, amssymb, amsthm
% Required environments: lemma, proposition, proof (standard with amsthm)
%
% The softmax attention entropy argument is kept as a paragraph-level remark
% rather than a formal result -- it's harder to formalize tightly and the
% effect direction is less clean than the normalization argument.
%
% The "compounding across layers" discussion is deliberately honest about the
% MLP/attention expansion counterforce rather than overclaiming. The net
% per-layer effect is compression minus expansion; we only prove the
% compression side.
%
% The unembedding argument (representation -> output diversity) relies on
% linearity of W_U. This gives an upper bound, which is the direction
% needed, but the bound could be loose if W_U has much lower rank than d.
%
% === END NOTES ===

\section{A Diversity-Reduction Theory of Activation Steering}
\label{sec:diversity-reduction}

We develop a formal account of how constant additive steering vectors reduce the diversity of a transformer's internal representations.
The core mechanism is the interaction between the fixed perturbation and the normalization layers (LayerNorm or RMSNorm) present in every transformer block: normalization projects out variation along the shifted mean and compresses the remaining variation by a factor that grows with the steering magnitude.
We derive the output covariance of a single normalization step under steering in closed form, characterize how effective rank and total variance depend on the steering magnitude, and discuss how the effect compounds across layers.

\subsection{Setup and Notation}

Let $x \in \mathbb{R}^d$ be a residual-stream vector drawn from a distribution $\mathcal{D}$ with mean $\mu = \mathbb{E}[x]$ and covariance $\Sigma_x = \mathbb{E}[(x - \mu)(x - \mu)^\top]$.
A steering intervention adds a fixed vector $s \in \mathbb{R}^d$ to the residual stream, producing $x' = x + s$.
This steered vector then passes through RMSNorm\footnote{%
  We analyze RMSNorm rather than full LayerNorm for clarity.
  LayerNorm additionally subtracts the per-vector mean, which projects onto a $(d{-}1)$-dimensional subspace orthogonal to the all-ones vector $\mathbf{1}$ \citep{brody2023expressivity}.
  This introduces a second rank-reducing projection but does not alter the core argument.
  In modern architectures (LLaMA, Qwen, etc.), RMSNorm has largely replaced LayerNorm.
}, which computes:
%
\begin{equation}
  z = \frac{x + s}{\|x + s\|_{\mathrm{RMS}}},
  \qquad
  \|v\|_{\mathrm{RMS}} = \sqrt{\frac{1}{d}\sum_{i=1}^{d} v_i^2}
  = \frac{\|v\|}{\sqrt{d}}.
  \label{eq:rmsnorm}
\end{equation}
%
Since $\|v\|_{\mathrm{RMS}} = \|v\| / \sqrt{d}$, RMSNorm is equivalent to projection onto the sphere of radius $\sqrt{d}$, followed by element-wise multiplication by learnable scale parameters $\gamma$.
Because $\gamma$ is fixed after training and does not depend on $x$, it acts as a constant linear transformation that does not affect rank.
We therefore analyze the normalized vector up to the constant $\sqrt{d}$:
%
\begin{equation}
  z = g(x + s), \qquad g(v) \coloneqq \frac{v}{\|v\|}.
  \label{eq:sphere-projection}
\end{equation}

\subsection{First-Order Expansion of the Normalization Map}

The function $g\colon \mathbb{R}^d \setminus \{0\} \to S^{d-1}$ is nonlinear, so we linearize it around the shifted mean $m \coloneqq \mu + s$ to obtain a tractable expression for the output covariance.
Write $x + s = m + \delta$ where $\delta = x - \mu$ satisfies $\mathbb{E}[\delta] = 0$ and $\mathbb{E}[\delta\delta^\top] = \Sigma_x$.

\begin{lemma}[Jacobian of the sphere projection]
\label{lem:jacobian}
The Jacobian of $g(v) = v / \|v\|$ is
\begin{equation}
  J_g(v) = \frac{1}{\|v\|}\left(I - \hat{v}\hat{v}^\top\right),
  \label{eq:jacobian}
\end{equation}
where $\hat{v} = v / \|v\|$ is the unit vector along $v$.
\end{lemma}

\begin{proof}
We compute the $(i,j)$-entry of $J_g$ directly.
Write $g_i(v) = v_i / h(v)$ where $h(v) = \|v\| = (v^\top v)^{1/2}$.
The partial derivative of $h$ with respect to $v_j$ is:
%
\begin{equation}
  \frac{\partial h}{\partial v_j}
  = \frac{1}{2}\left(\sum_k v_k^2\right)^{-1/2} \cdot 2v_j
  = \frac{v_j}{\|v\|}.
\end{equation}
%
By the quotient rule:
%
\begin{equation}
  \frac{\partial g_i}{\partial v_j}
  = \frac{\delta_{ij} \cdot \|v\| - v_i \cdot \frac{v_j}{\|v\|}}{\|v\|^2}
  = \frac{\delta_{ij}}{\|v\|} - \frac{v_i v_j}{\|v\|^3}.
\end{equation}
%
In matrix form, the first term is $\frac{1}{\|v\|}I$ and the second is $\frac{1}{\|v\|^3}vv^\top = \frac{1}{\|v\|}\hat{v}\hat{v}^\top$, giving \eqref{eq:jacobian}.
\end{proof}

Applying the first-order Taylor expansion $g(m + \delta) \approx g(m) + J_g(m)\,\delta$:
%
\begin{equation}
  z \approx \hat{m} + \frac{1}{\|m\|}\,P_{\perp m}\,\delta,
  \label{eq:taylor}
\end{equation}
%
where $P_{\perp m} \coloneqq I - \hat{m}\hat{m}^\top$ is the orthogonal projector onto the subspace perpendicular to $m$.

\paragraph{Validity of the approximation.}
The expansion \eqref{eq:taylor} is accurate when $\|\delta\| \ll \|m\|$, i.e., the input variation is small relative to the shifted mean.
This condition is increasingly well-satisfied as the steering magnitude $\|s\|$ grows (since $\|m\| \geq \|s\| - \|\mu\|$\footnote{%
  This is the reverse triangle inequality:
  $\|a + b\| \geq \bigl|\|a\|-\|b\|\bigr| \geq \|a\|-\|b\|$.
  One-line proof: write $s = m + (-\mu)$ and apply the ordinary triangle inequality,
  $\|s\| \leq \|m\| + \|\mu\|$, then rearrange.
  Geometrically, $m = \mu + s$ lies in a ball of radius $\|\mu\|$ centered at $s$, so
  the closest $m$ can come to the origin is $\|s\| - \|\mu\|$ (achieved when $\mu \parallel -s$).
  The bound is only informative once $\|s\| > \|\mu\|$; otherwise $\|s\|-\|\mu\|$ is negative
  and the inequality is vacuous.
  We use this step in several places below (e.g.\ Section~\ref{sec:covariance}, and in the
  geometric-bound proofs of Section~\ref{sec:geometric-bound}).%
}), which is precisely the regime where the diversity-reduction effect is strongest.
The approximation breaks down when $s \approx -\mu$ so that $\|m\| \approx 0$; this anti-pole regime is analyzed in Section~\ref{sec:conditions}.
Second-order corrections are of order $O(\|\delta\|^2 / \|m\|^2)$.

\subsection{Covariance Under Steering}
\label{sec:covariance}

From \eqref{eq:taylor}, the output has mean $\mathbb{E}[z] \approx \hat{m}$ and fluctuation $z - \hat{m} \approx \frac{1}{\|m\|}P_{\perp m}\,\delta$.
The output covariance is therefore, using the fact that $P_{\perp m}$ is
symmetric\footnote{%
  $P_{\perp m}^\top = (I - \hat{m}\hat{m}^\top)^\top = I - \hat{m}\hat{m}^\top = P_{\perp m}$,
  because $\hat{m}\hat{m}^\top$ is an outer product of a vector with itself.%
}:
%
\begin{equation}
  \Sigma_z
  \;=\; \mathbb{E}\!\left[(z - \hat{m})(z - \hat{m})^\top\right]
  \;\approx\; \frac{1}{\|m\|^2}\,P_{\perp m}\,\Sigma_x\,P_{\perp m}^\top
  \;=\; \frac{1}{\|m\|^2}\,P_{\perp m}\,\Sigma_x\,P_{\perp m}.
  \label{eq:output-cov-simplified}
\end{equation}

Two effects are visible in \eqref{eq:output-cov-simplified}:
\begin{enumerate}
  \item \textbf{Rank reduction via projection.}
    The projector $P_{\perp m}$ has rank $d - 1$: it annihilates all variation along $\hat{m}$.
    Geometrically, projecting onto the unit sphere eliminates radial variation---any component of $\delta$ along $m$ changes only the norm $\|m + \delta\|$, not the direction $\widehat{m + \delta}$.
    The direction that is annihilated, $\hat{m} = \widehat{\mu + s}$, depends on $s$; different steering vectors eliminate different dimensions of variation.

  \item \textbf{Variance shrinkage via rescaling.}
    All surviving variation is scaled by $1/\|m\|^2$.
    As $\|s\|$ grows, $\|m\| = \|\mu + s\|$ grows (at least as fast as $\|s\| - \|\mu\|$), compressing the total variance toward zero.
    This is the formalization of the informal observation by \citet{turntrout2023steering} that when the steering coefficient is large, the steering vector ``comprises most of the residual stream content'' after LayerNorm.
\end{enumerate}

\begin{proposition}[Covariance under steered RMSNorm]
\label{prop:diversity-reduction}
Let $\mathcal{X}$ be a distribution over $\mathbb{R}^d$ with covariance $\Sigma_x$, let $s \in \mathbb{R}^d$ be a fixed steering vector, and let $m = \mathbb{E}[x] + s$.
Under the first-order approximation \eqref{eq:taylor}, the covariance $\Sigma_z$ of the RMSNorm output satisfies:
%
\begin{enumerate}
  \item \textup{(Rank reduction)}
    $\mathrm{rank}(\Sigma_z) \leq \min\!\big(\mathrm{rank}(\Sigma_x),\; d - 1\big)$.
    If $\Sigma_x$ has nonzero variance along $\hat{m}$ (i.e., $\hat{m}^\top \Sigma_x \hat{m} > 0$), then $\mathrm{rank}(\Sigma_z) \leq \mathrm{rank}(\Sigma_x) - 1$.

  \item \textup{(Total variance reduction)}
    $\mathrm{tr}(\Sigma_z) = \frac{1}{\|m\|^2}\big(\mathrm{tr}(\Sigma_x) - \hat{m}^\top \Sigma_x\, \hat{m}\big)$.
\end{enumerate}
\end{proposition}

\begin{proof}
For the rank bound: $\Sigma_z = \frac{1}{\|m\|^2} P_{\perp m}\Sigma_x P_{\perp m}$, so $\mathrm{rank}(\Sigma_z) = \mathrm{rank}(P_{\perp m}\Sigma_x P_{\perp m})$.
Since $P_{\perp m}$ has rank $d-1$, the product has rank at most $\min(\mathrm{rank}(\Sigma_x), d-1)$.
If $\hat{m}^\top \Sigma_x \hat{m} > 0$, then $\Sigma_x$ has a component in the nullspace of $P_{\perp m}$, and the rank drops by at least one.

For the total variance: using the cyclic property of the trace, $\mathrm{tr}(ABC) = \mathrm{tr}(CAB)$, and the idempotence of $P_{\perp m}$:
%
\begin{align}
  \mathrm{tr}(\Sigma_z)
  &= \frac{1}{\|m\|^2}\,\mathrm{tr}(P_{\perp m}\,\Sigma_x\,P_{\perp m})
  \label{eq:trace-step1}\\
  &= \frac{1}{\|m\|^2}\,\mathrm{tr}(P_{\perp m}^2\,\Sigma_x)
  \label{eq:trace-step2}\\
  &= \frac{1}{\|m\|^2}\,\mathrm{tr}(P_{\perp m}\,\Sigma_x)
  \label{eq:trace-step3}\\
  &= \frac{1}{\|m\|^2}\,\mathrm{tr}\!\big((I - \hat{m}\hat{m}^\top)\Sigma_x\big)
  \label{eq:trace-step4}\\
  &= \frac{1}{\|m\|^2}\big(\mathrm{tr}(\Sigma_x) - \mathrm{tr}(\hat{m}\hat{m}^\top\Sigma_x)\big)
  \label{eq:trace-step5}\\
  &= \frac{1}{\|m\|^2}\big(\mathrm{tr}(\Sigma_x) - \hat{m}^\top\Sigma_x\,\hat{m}\big).
  \label{eq:trace-step6}
\end{align}
%
The step from \eqref{eq:trace-step1} to \eqref{eq:trace-step2} uses the cyclic property to move the rightmost $P_{\perp m}$ to the left.
The step from \eqref{eq:trace-step2} to \eqref{eq:trace-step3} uses idempotence: $P_{\perp m}^2 = P_{\perp m}$, which holds because
$(I - \hat{m}\hat{m}^\top)^2 = I - 2\hat{m}\hat{m}^\top + \hat{m}(\hat{m}^\top\hat{m})\hat{m}^\top = I - \hat{m}\hat{m}^\top$,
where $\hat{m}^\top\hat{m} = 1$ since $\hat{m}$ is a unit vector.
The step from \eqref{eq:trace-step5} to \eqref{eq:trace-step6} uses the cyclic property to write $\mathrm{tr}(\hat{m}\hat{m}^\top\Sigma_x) = \mathrm{tr}(\hat{m}^\top\Sigma_x\hat{m}) = \hat{m}^\top\Sigma_x\hat{m}$, where the last equality holds because $\hat{m}^\top\Sigma_x\hat{m}$ is a scalar.
\end{proof}

\paragraph{Interpretation.}
The numerator $\mathrm{tr}(\Sigma_x) - \hat{m}^\top\Sigma_x\hat{m}$ is the original total variance minus the variance along the direction $\hat{m} = \widehat{\mu + s}$.
For generic $s$, this quantity is strictly less than $\mathrm{tr}(\Sigma_x)$.
The denominator $\|m\|^2 = \|\mu + s\|^2$ grows quadratically with the steering magnitude.
In the strong-steering regime $\|s\| \gg \|\mu\|$:
%
\begin{equation}
  \mathrm{tr}(\Sigma_z) \approx \frac{\mathrm{tr}(\Sigma_x) - \hat{s}^\top\Sigma_x\,\hat{s}}{\|s\|^2},
\end{equation}
%
which vanishes as $\|s\| \to \infty$.

\subsection{Secondary Mechanisms}

\paragraph{Softmax attention entropy.}
Steering shifts the query vectors by a constant, modifying the attention logits $q^\top k_i$.
Because softmax is a concave function on the probability simplex with respect to uniform perturbations of the logits, a constant shift toward particular key positions reduces the entropy of the attention distribution.
Lower attention entropy means fewer positions contribute to the output, further reducing the diversity of information flowing through the layer.
We note this as a secondary contributor that is harder to formalize tightly.

\paragraph{Compounding across layers.}
In a pre-norm transformer, each block applies the sequence:
$\mathrm{Norm} \to \mathrm{Attn} \to \mathrm{Add} \to \mathrm{Norm} \to \mathrm{MLP} \to \mathrm{Add}$.
The steering vector $s$, once added to the residual stream, persists via the residual connections.
At each subsequent normalization step, the projection and shrinkage from Proposition~\ref{prop:diversity-reduction} apply with the current shifted mean at that layer.
In principle, effective rank can drop by up to one at each normalization step.
However, the MLP and attention sub-layers are nonlinear maps that can expand variation (e.g., a ReLU/GeLU activation can map low-variance inputs to higher-variance outputs if the inputs straddle the nonlinearity).
The net effect at each layer is therefore normalization-induced compression minus sub-layer-induced expansion.
\citet{sinii2025small} provide direct empirical evidence for the normalization bottleneck: they find that steering vectors placed before the input LayerNorm of a block lose effectiveness, and that placing the vector \emph{after} the LayerNorm restores performance---indicating that the LayerNorm step specifically attenuates the steering signal.

\paragraph{From representation diversity to output diversity.}
The unembedding layer is a linear map $W_U \in \mathbb{R}^{V \times d}$ (where $V$ is vocabulary size) that converts residual-stream vectors to logits.
The covariance of the logit distribution is $W_U \Sigma_z W_U^\top$, whose rank is at most $\mathrm{rank}(\Sigma_z)$.
Since rank cannot increase through linear maps, representation-level diversity reduction is an \emph{upper bound} on output-level diversity, which is the direction needed for our claims.

\subsection{Relationship to Prior Work}

\citet{dong2021attention} showed that pure self-attention without skip connections suffers from doubly exponential rank collapse with depth.
\citet{wu2024role} extended this analysis to include LayerNorm, showing that LayerNorm can either prevent or accelerate collapse depending on the value matrices---importantly establishing that the rank dynamics of self-attention with LayerNorm are richer than previously supposed.
Our analysis differs from this line of work in that we study rank reduction caused not by depth but by an \emph{external additive perturbation} interacting with normalization.
The depth-collapse results concern the unsteered model's intrinsic dynamics; our result concerns a controlled intervention.

\citet{turntrout2023steering} informally observed that GPT-2's LayerNorm renormalizes the residual stream to $\sqrt{d_{\mathrm{model}}}$, so that a large steering coefficient causes the steering vector to dominate the post-norm content.
Our Proposition~\ref{prop:diversity-reduction} formalizes this observation: the ``domination'' is precisely the variance shrinkage by $1/\|m\|^2$, and the ``most of the residual stream content'' corresponds to the projection $P_{\perp m}$ annihilating variation along the steering direction.
The companion post on residual stream norm growth \citep{turntrout2023norms} establishes that model weights compensate for LayerNorm's normalization by producing exponentially growing outputs---a mechanism that steering vectors, being fixed and not adapted to the model's learned growth schedule, do not participate in.

Most directly, \citet{sinii2025small} empirically demonstrate that the input LayerNorm is the specific bottleneck limiting steering vector effectiveness between adjacent layers, and that bypassing it recovers performance.
They frame this as an engineering obstacle; we reframe it as the \emph{primary architectural mechanism} through which steering compresses representational diversity.

On the output side, \citet{zhang2025verbalized} attribute post-alignment mode collapse to typicality bias in preference data, and \citet{omahony2024attributing} attribute it to the mode-seeking nature of reverse KL in RLHF/DPO.
These are complementary explanations at the level of training signal; our account provides a mechanistic explanation at the level of architecture for why inference-time steering interventions produce analogous diversity reduction without any training.
